% ==================================================================
% id:       der-001
% topic:    נגזרות
% subtopic: הגדרת הנגזרת וגזירות
% solution: full
% ==================================================================
\textbf{עבור כל אחת מהפונקציות הבאות, חשבו את הנגזרת שלה, בכל נקודה בתחום ההגדרה,  ע"פ הגדרת הנגזרת, ללא שימוש בכללי חישוב נגזרת:}
\bE\bMs{2}
\item $\ds{f(x)=\frac{1}{x^3}}$
\hint{הביאו את המונה של מנת ההפרשים למכנה משותף ופתחו את $(x+h)^3$, כך שיהיה אפשר לצמצם ב-$h$.}
\item $\ds{f(x)=\cos(5x)}$
\hint{השתמשו בזהות להפרש קוסינוסים כדי להגיע לגבול המופלא $\frac{\sin t}{t}\to1$.}
\item $\ds{f(x)=\ln\lr{1 + x^2}}$
\hint{כתבו את הפרש הלוגריתמים כלוגריתם של מנה מהצורה $1+t$ עם $t\to0$, והשתמשו בגבול $\frac{\ln(1+t)}{t}\to1$.}
\item $\ds{f(x)=e^{\sin x}}$
\hint{הוציאו את $e^{\sin x}$ כגורם משותף והשתמשו בגבול $\frac{e^t-1}{t}\to1$, כאשר $t$ הוא הפרש ערכי הסינוס.}
\eMs
\eE
\ifsol
\begin{solbox}
\bE
\item
מתקיים
\begin{align*}
    f'(x)&=\ \lim_{h\to0}\frac{f(x+h)-f(x)}{h}\\
    &=\ \lim_{h\to0}\frac{\tfrac{1}{(x+h)^3}-\tfrac{1}{x^3}}{h},
\end{align*}
לפני שנחשב את הגבול, נפשט:
\begin{align*}
    \frac{\tfrac{1}{(x+h)^3}-\tfrac{1}{x^3}}{h}&=\ \frac{x^3-(x+h)^3}{hx^3(x+h)^3}\\
    &=\ \frac{-3x^2h-3xh^2-h^3}{hx^3(x+h)^3}\\
    &=\ \frac{-3x^2-3xh-h^2}{x^3(x+h)^3},
\end{align*}
מכאן
\begin{align*}
    f'(x)&=\ \lim_{h\to0}\frac{-3x^2-3xh-h^2}{x^3(x+h)^3}\\
    &=\ \frac{-3x^2-3x\cdot0-0^2}{x^3(x+0)^3}= -\frac{3}{x^4}.
\end{align*}

\item
מתקיים
\begin{align*}
    f'(x_0)&=\ \lim_{x\to x_0}\frac{f(x)-f(x_0)}{x-x_0}\\
    &=\ \lim_{x\to x_0}\frac{\cos(5x)-\cos(5x_0)}{x-x_0},
\end{align*}
נזכר בזהות
$$.\cos\alpha-\cos\beta=-2\sin(\tfrac{\alpha-\beta}{2})\sin(\tfrac{\alpha+\beta}{2})$$
נפשט
\begin{align*}
    \frac{\cos(5x)-\cos(5x_0)}{x-x_0}&=\ -2\frac{\sin(\tfrac{5x-5x_0}{2})\sin(\tfrac{5x+5x_0}{2})}{x-x_0}\\
    &=\ -5\cdot\frac{\sin(\tfrac{5(x-x_0)}{2})}{\tfrac{5(x-x_0)}{2}}\cdot\sin(\tfrac{5x+5x_0}{2}),
\end{align*}
מכאן
\begin{align*}
    f'(x_0)&=\ -5\lim_{x\to x_0}\frac{\sin(\tfrac{5(x-x_0)}{2})}{\tfrac{5(x-x_0)}{2}}\cdot\sin(\tfrac{5x+5x_0}{2})\\
    &\overset{(\ast)}{=}\ -5\cdot1\cdot\sin(\tfrac{10x_0}{2})\\
    &=\ -5\sin(5x_0).
\end{align*}
$(\ast)$
יש כאן שימוש בגבול מופלא
$\displaystyle\lim_{t\to0}\frac{\sin t}{t}=1$.
\item
מתקיים
\begin{align*}
    f'(x)&=\ \lim_{h\to0}\frac{f(x+h)-f(x)}{h}\\
    &=\ \lim_{h\to0}\frac{\ln(1+(x+h)^2)-\ln(1+x^2)}{h},
\end{align*}
לפני שנחשב את הגבול, נפשט:
\begin{align*}
    \frac{\ln(1+(x+h)^2)-\ln(1+x^2)}{h}&=\ \frac{\ln\left(\tfrac{1+(x+h)^2}{1+x^2}\right)}{h}\\
    &=\ \frac{\ln\left(1+\tfrac{2xh+h^2}{1+x^2}\right)}{h},
\end{align*}
נשים לב כי
$$,\lim_{h\to 0}\frac{2xh+h^2}{1+x^2}=0$$
ונמשיך
\begin{align*}
    \frac{\ln\left(1+\tfrac{2xh+h^2}{1+x^2}\right)}{h}&=\ \frac{\ln\left(1+\tfrac{2xh+h^2}{1+x^2}\right)}{h{\color{red}\cdot\frac{2xh+h^2}{1+x^2}}}{\color{red}\cdot\frac{2xh+h^2}{1+x^2}}\\
    &=\ \frac{\ln\left(1+\tfrac{2xh+h^2}{1+x^2}\right)}{{\frac{2xh+h^2}{1+x^2}}}{\cdot\frac{2x+h}{1+x^2}}\\
\end{align*}
מכאן
\begin{align*}
    f'(x)&=\ \lim_{h\to0}\frac{\ln\left(1+\tfrac{2xh+h^2}{1+x^2}\right)}{{\frac{2xh+h^2}{1+x^2}}}{\cdot\frac{2x+h}{1+x^2}}\\
    &\overset{(\ast)}{=}\ 1\cdot\frac{2x+0}{1+x^2}\\
    &=\ \frac{2x}{1+x^2}.
\end{align*}

$(\ast)$
יש כאן שימוש בגבול מופלא
$\displaystyle\lim_{t\to0}\frac{\ln(1+t)}{t}=1$.
\item
מתקיים
\begin{align*}
    f'(x)&=\ \lim_{h\to0}\frac{f(x+h)-f(x)}{h}\\
    &=\ \lim_{h\to0}\frac{e^{\sin(x+h)}-e^{\sin x}}{h},
\end{align*}
לפני שנחשב את הגבול, נפשט:
\begin{align*}
    \frac{e^{\sin(x+h)}-e^{\sin x}}{h}&=\ e^{\sin x}\cdot\frac{e^{\sin(x+h)-\sin x}-1}{h},
\end{align*}
נשים לב כי
$$,\lim_{h\to0}(\sin(x+h)-\sin x)=0$$
נמשיך לפשט
\begin{align*}
    e^{\sin x}\cdot\frac{e^{\sin(x+h)-\sin x}-1}{h}&=\ e^{\sin x}\cdot\frac{e^{\sin(x+h)-\sin x}-1}{{\color{red}(\sin(x+h)-\sin x)\cdot}h}{\color{red}\cdot(\sin(x+h)-\sin x)}\\
    &=\ e^{\sin x}\cdot\frac{e^{\sin(x+h)-\sin x}-1}{{\color{red}(\sin(x+h)-\sin x)\cdot}}\cdot\frac{{\sin(x+h)-\sin x}}{h}
\end{align*}
מכאן
\begin{align*}
    f'(x)&=\ e^{\sin x}\cdot\lim_{h\to0}\frac{e^{\sin(x+h)-\sin x}-1}{{\color{red}(\sin(x+h)-\sin x)\cdot}}\cdot\frac{{\sin(x+h)-\sin x}}{h}\\
    &\overset{(\ast)}{=}\ e^{\sin x}\cdot1\cdot\cos x\\
    &=\ e^{\sin x}\cdot\cos x.
\end{align*}
 $(\ast)$
 יש כאן שימוש בגבול מופלא
 $\lim_{t\to0}(1+t)^{\frac{1}{t}}=e$.
 כמוכן, שימו לב כי
 $$\displaystyle\lim_{h\to 0}\frac{{\sin(x+h)-\sin x}}{h}=\sin'(x)=\cos x$$
\eE
\end{solbox}
\fi
